\chapter{Parent Inference and Controlled Memory}
\label{ch:inference-memory}

A normalized parent law does not by itself specify parent-sensitive inference
or a closed dynamics on a smaller observation. This chapter supplies a finite
inference mechanism, identifies the preparation under which its full parent
is Markov, and retains the exact hidden history when a further projection
has memory. Finite-memory approximation then has its own normalized path
law and an exact information loss. The transition index is an auxiliary
inference index; it is not physical time. \status{DEFINITION}

\section{A fixed normalized target and a parent-sensitive sweep}
\label{sec:rg-memory-parent-mechanism}

Let \(\mathsf I\) be a nonempty finite incidence set with projections
\(z:\mathsf I\to\mathsf Z\) and \(h:\mathsf I\to\mathsf H\), where
\(h\) retains the complete decorated parent. Adopt positive generative
weights \(g_i\), normalized record conditionals \(\eta_i(e)\), and
normalized observation likelihoods \(\ell_i(o)\), all independent of
recognition. At an observation for which all retained \(\ell_i(o)>0\),
define
\begin{equation}
 Z_{\mathrm{gen}}=\sum_i g_i,\qquad Z_o=\sum_i g_i\ell_i(o),
 \qquad p_i=\frac{g_i\ell_i(o)}{Z_o}.
\label{eq:rg-memory-target}
\end{equation}
The complete generative law is
\(g_i\eta_i(e)\ell_i(o)/Z_{\mathrm{gen}}\). The local symbol \(p\)
denotes this finite posterior, not a new agent. Positivity and the specified
incidence support are adopted hypotheses for the finite results below.
\status{HYPOTHESIS}

For recognition \(q_i\eta_i(e)\), direct substitution in the exact ELBO
gives the established identity
\begin{equation}
 \mathcal F(q;o)=\KL(q\Vert p)-\log(Z_o/Z_{\mathrm{gen}}).
\label{eq:rg-memory-fixed-vfe}
\end{equation}
A different conditional record recognition retains its averaged conditional
KL. Thus matching the displayed conditionals is a specialization of the
complete joint-law objective, not removal of its dependence terms.
\status{ESTABLISHED}

At fixed \(z\), declare an undirected graph of allowed parent changes and
a symmetric proposal \(a_{ij}=a_{ji}\geq0\) whose off-diagonal row sum
is at most \(1/2\). Its Metropolis kernel, the parent projection, and its
posterior conditional lift are
\begin{equation}
\begin{aligned}
 M_{ij}&=a_{ij}\min(1,p_j/p_i)\quad(i\ne j),
 &M_{ii}&=1-\sum_{j\ne i}M_{ij},\\
 C_{iH}&=\mathbf1_{h(i)=H},
 &J_{Hi}&=\frac{p_i\mathbf1_{h(i)=H}}{p_H(H)},
 \qquad p_H=pC,\\
 D&=CJ,&T&=MD.
\end{aligned}
\label{eq:rg-memory-sweep}
\end{equation}
Only nonempty parent fibers are retained. The step \(D\) refreshes the
fine conditional law at the current complete parent. It is nontrivially
parent-sensitive precisely when those fine conditional laws differ. The
proposal and update order are declared mechanism choices.
\status{DEFINITION}

The complete marked version of a row kernel \(K\) is
\(\widetilde K((i,e),(j,e'))=K_{ij}\eta_j(e')\). Summing the terminal
record proves normalization, and summing the initial record proves that
\(p_i\eta_i(e)\) is invariant whenever \(pK=p\). The matched family
\(q_i\eta_i(e)\) is preserved. These are established lift identities;
unmatched conditional recognition retains its own information terms.
\status{ESTABLISHED}

One finite assignment graph uses interval partitions of four ordered children
and one binary root mark per block. A move flips one existing mark, merges
adjacent blocks retaining the left mark, or reverses such a merge while
choosing the exposed right mark. There are
\(\sum_{r=1}^4\binom3{r-1}2^r=54\) assignments and seven neighbors per
assignment: for \(r\) blocks the counts are \(r\), \(r-1\), and
\(2(4-r)\). Merge-to-one and split-to-target paths show connectivity.
Consequently proposal \(1/14\) per edge satisfies the preceding conditions.
This is an established finite realization of the move graph, not an intrinsic
choice of partition or parent. \status{ESTABLISHED}

\theoremheading{Stationarity, exact descent, and finite-component convergence}{thm:rg-memory-inference-descent}
The kernels \(M,D,T\) are normalized and preserve \(p\). Let a component
mean a connected component of the union of their positive off-diagonal
edges. For every initial recognition \(q\), \(T\) preserves component
masses and decreases \cref{eq:rg-memory-fixed-vfe}, strictly unless
\(q(i)=q(\mathsf B)p(i\mid\mathsf B)\) on each component \(\mathsf B\).
On each fixed component it converges geometrically to that conditional target.
These are established finite-instance results; they give no uniform rate as
state counts grow or positive weights tend to zero.
\status{ESTABLISHED}

\paragraph{Proof.}
For an edge, \(p_iM_{ij}=a_{ij}\min(p_i,p_j)=p_jM_{ji}\), so \(M\)
is reversible. Also \(JC=I\) and \(pD=p\), while
\(p_iD_{ij}=p_ip_j/p_H(h(i))\) on a common parent fiber, so \(D\) is
reversible. Normalization follows from their definitions, and stationarity
survives composition. For any stationary finite kernel \(K\), put
\(q'=qK\) and \(K^*_{ji}=p_iK_{ij}/p_j\). The two input-output joint
laws \(q_iK_{ij}\) and \(p_iK_{ij}\) have KL \(\KL(q\Vert p)\).
Factoring by their output marginal proves
\begin{equation}
 \KL(q\Vert p)-\KL(qK\Vert p)
 =\sum_{j:q'_j>0}q'_j
   \KL\bigl(q(i\mid j)\Vert K^*_{j\cdot}\bigr).
\label{eq:rg-memory-reverse-conditional-loss}
\end{equation}
Every summand is nonnegative. With positive diagonal, equality requires the
ratio \(q_i/p_i\) to be constant across every positive edge: the output
conditional equality identifies that ratio at \(i\) and \(j\), and the
holding edge identifies it with the output ratio. Equality for \(M\)
therefore makes \(q\) stationary under \(M\); equality for the subsequent
\(D\) additionally makes the ratio constant on parent fibers. This is
exactly constancy on the stated union components.

The positive holding edges of both factors ensure that every \(M\) or
\(D\) edge is a positive \(T\) edge. On a component with \(n\) states,
a shortest path has length at most \(n-1\); positive holding extends it to
one common length \(m\geq\max(1,n-1)\). Hence every entry of
\(T^m\) on that component is positive. Writing its normalized target as
\(\pi\), let \(\varepsilon=\min_{ij}(T^m)_{ij}/\pi_j>0\).
Then \(T^m=\varepsilon\mathbf1\pi+(1-\varepsilon)R\) with a
stationary stochastic \(R\), unless \(\varepsilon=1\), when convergence
is exact. Convexity and data processing give
\(\KL(qT^m\Vert\pi)\leq(1-\varepsilon)\KL(q\Vert\pi)\), and
the same mixture contracts total variation. Iteration proves convergence.
The original full-model VFE on this component includes
\(-\log p(\mathsf B)-\log(Z_o/Z_{\mathrm{gen}})\); neither constant
is deleted by the conditional dynamics. \(\square\)

\theoremheading{The prepared complete parent is Markov}{thm:rg-memory-prepared-parent}
Define \(L=JMC\). Then \(L\) is reversible under \(p_H\), and
\begin{equation}
 JT=LJ,\qquad (\nu J)T^n=\nu L^nJ.
\label{eq:rg-memory-parent-semiconjugacy}
\end{equation}
For a start \(q_0=\nu J\), the complete parent history has transition
\(L\). For an arbitrary start, this statement applies after the first
completed \(D\) stage. It is an established prepared-history statement,
not universal closure of every parent projection at the initial instant.
\status{ESTABLISHED}

\paragraph{Proof.}
Use \(T=MCJ\) and \(JC=I\) to obtain \(JT=JMCJ=LJ\), then
induct. Moreover
\(p_H(H)L_{HH'}=\sum_{i:h(i)=H,j:h(j)=H'}p_iM_{ij}\), which is
symmetric by reversibility of \(M\). The final conditional refresh makes
the fine law given the entire preceding parent history equal to
\(J_{H\cdot}\), depending only on its last parent. The next parent row
is consequently \(J_{H\cdot}MC=L_{H\cdot}\). Induction over histories
proves the Markov assertion, including its preparation condition.
\(\square\)

\section{What persistence and projection actually imply}
\label{sec:rg-memory-persistence}

For a finite irreducible component with stationary law \(\pi\), choose a
nonempty proper event \(A\), such as a specified retained block event.
Let \(K_A\) be the restriction of its transition to \(A\). For an
initial law \(\mu\) supported in \(A\), the event of remaining inside
for \(n\) further transitions has probability
\begin{equation}
 s_n=\mu K_A^n\mathbf1,\qquad
 \E_\mu\tau_{A^c}=\sum_{n\geq0}s_n
                 =\mu(I-K_A)^{-1}\mathbf1.
\label{eq:rg-memory-killed-survival}
\end{equation}
These are established formulas for unbroken survival with lost mass retained
as lost mass. \status{ESTABLISHED}

\paragraph{Proof.}
Multiplication by the restricted matrix sums exactly the paths staying in
\(A\). For a positive integer-valued exit time, summing its tail gives
its mean by summing nonnegative indicator functions. Every state in \(A\)
has a finite positive exit path. Finiteness supplies a common number of steps
and a positive lower exit probability, so the killed tail is geometric.
The convergent matrix series is the inverse displayed above. \(\square\)

\theoremheading{Stationary episode means and the clock control}{thm:rg-memory-episode-means}
Let \(J_A=\sum_{i\notin A,j\in A}\pi_iK_{ij}\). Under the stationary
entrance preparation, the inside and outside episode means are
\begin{equation}
 \E\tau_{A^c}=\frac{\pi(A)}{J_A},\qquad
 \E\tau_A=\frac{\pi(A^c)}{J_A}.
\label{eq:rg-memory-stationary-episodes}
\end{equation}
For whole-clock waiting \(K_\lambda=(1-\lambda)I+\lambda K\),
\(0<\lambda\leq1\), every finite fixed-preparation hitting mean is
multiplied by \(1/\lambda\). These established identities separate
stationary occupancy, crossing frequency, and survival shape.
\status{ESTABLISHED}

\paragraph{Proof.}
The unnormalized entrance row on \(A\) is
\(r=\pi_{A^c}K_{A^c,A}\). Stationarity gives
\(\pi_A=\pi_AK_A+r\), so
\(r(I-K_A)^{-1}=\pi_A\). Divide by \(J_A=r\mathbf1\) and use
\cref{eq:rg-memory-killed-survival}. Stationarity equates entrance and
exit flux, giving the outside formula as well. Finally
\(I-(K_\lambda)_A=\lambda(I-K_A)\), so its inverse rescales every
finite mean by \(1/\lambda\). Equal stationary occupancy and flux
therefore force equal stationary episode means but need not give equal
survival distributions. \(\square\)

The finite mechanism does not establish universal persistent hierarchy,
autonomous agency, or a preferred clock. Establishing a persistence advantage
requires a stated event, preparation, and comparator after accounting for
occupancy and waiting. Higher-level lineage and geometry-changing dynamics
remain separate open construction problems. \status{OPEN}

\theoremheading{The exact criterion for a universal projected chain}{thm:rg-memory-lumpability}
For a finite hidden kernel \(K\) and deterministic observation
\(f:\mathsf S\to\mathsf A\), one observed transition kernel works for
every initial law exactly when
\begin{equation}
 \sum_{j:f(j)=b}K_{ij}
 =\sum_{j:f(j)=b}K_{i'j}
 \quad\text{whenever }f(i)=f(i'),\quad b\in\mathsf A.
\label{eq:rg-memory-lumpability}
\end{equation}
This condition is necessary and sufficient for the complete supported rows. \status{ESTABLISHED}

\paragraph{Proof.}
Point-mass starts at \(i,i'\) prove necessity. Under the condition, every
hidden posterior given an observed history ending at \(a\) is supported on
\(f^{-1}(a)\), and all rows there give the same next-observation law.
Averaging therefore gives that law independently of earlier observations.
\(\square\)

Even a stationary observation can have memory. On the deterministic
three-state cycle with uniform stationary law and labels \((0,0,1)\),
conditioning on observed histories \((0,0)\) and \((1,0)\) gives next
symbols \(1\) and \(0\), respectively. Both histories are positive and
have the same final symbol. This is an established failure of first-order
Markov closure under that fixed stationary preparation.
\status{ESTABLISHED}

\section{The complete hidden history and its conditional state}
\label{sec:rg-memory-hidden-history}

For the rest of this chapter let \(Z_n\) be a finite hidden state with
initial law \(\nu\) and transition \(K\), and let
\(Y_n=f(Z_n)\) be its finite observation. These are local history symbols,
not the structural datum \(X\) or the original sensor observation \(o\).
Let \(E_a\) be the diagonal indicator of \(f^{-1}(a)\). The declared
joint finite path law and its observed marginal are
\begin{equation}
 \nu(z_0)\prod_{t=0}^{N-2}K_{z_tz_{t+1}}
       \prod_{t=0}^{N-1}\mathbf1_{f(z_t)=y_t},\qquad
 \mathbb P_N(y)=\nu E_{y_0}K E_{y_1}\cdots K E_{y_{N-1}}\mathbf1.
\label{eq:rg-memory-hidden-path}
\end{equation}
This auxiliary model is fixed before its path recognition is varied.
\status{DEFINITION}

\theoremheading{Exact filtering without a projected Markov assumption}{thm:rg-memory-filter}
For a positive word \(w=(y_0,\ldots,y_n)\), put
\begin{equation}
 a_w=\nu E_{y_0}K E_{y_1}\cdots K E_{y_n},\qquad
 \widehat\nu_w=\frac{a_w}{a_w\mathbf1},\qquad
 \mathbb P(Y_{n+1}=b\mid w)=\widehat\nu_wK E_b\mathbf1.
\label{eq:rg-memory-filter}
\end{equation}
These formulas define the exact terminal posterior and next prediction.
Every zero-mass word has zero messages for all extensions and undefined
conditional readouts. The terminal posterior determines every future suffix,
but its possible values need not form a finite set. The finite formulas and
support statements are established. \status{ESTABLISHED}

\paragraph{Proof.}
Summing the hidden path law while fixing its final state gives \(a_w\).
Successive summation of \(E_b\) to \(I\), followed by \(K\mathbf1
=\mathbf1\), proves normalization and prefix consistency. Division by
positive word mass gives the posterior; a further matrix product gives any
future suffix. Nonnegative zero-mass messages are identically zero and stay
zero on multiplication. A posterior is a probability vector on the hidden
set, so its sufficiency does not assert finiteness of its attainable subset.
\(\square\)

An auxiliary observed-word evidence and path VFE may be formed from
\cref{eq:rg-memory-hidden-path}. If \(\nu\) was chosen using a
recognition law of the original model, that choice does not turn this path
experiment into the original recognition-independent sensor model. The two
VFE identities have different evidences and must remain separately typed.
This identification is deliberately declined. \status{NOT-CLAIMED}

\section{Normalized finite-memory projection}
\label{sec:rg-memory-projection}

First take any fixed finite law \(\mathbb P_N\) on \(\mathsf A^N\),
without stationarity. For an integer \(k\geq0\) and position \(t\),
let \(c_t=y_{\max(0,t-k):t-1}\). On positive context mass define
\begin{equation}
 R_t^{(k)}(a\mid c)=\mathbb P_N(Y_t=a\mid c_t=c),\qquad
 \widehat{\mathbb P}_N^{(k)}(y)
    =\prod_{t=0}^{N-1}R_t^{(k)}(y_t\mid c_t).
\label{eq:rg-memory-projection-product}
\end{equation}
A zero prefix has zero extensions, so the product is defined recursively
without choosing a row on a null context. The hat marks a path approximation,
not a recognition law in the original statistical experiment.
\status{DEFINITION}

\theoremheading{Forward-KL optimality and its exact information defect}{thm:rg-memory-projection}
The product in \cref{eq:rg-memory-projection-product} is a normalized
law, has the same retained context-plus-symbol marginals as
\(\mathbb P_N\), and is positive on every \(\mathbb P_N\)-positive
word. Among all normalized causal laws with time-dependent \(k\)-context
rows it uniquely minimizes the forward path KL. Its minimum is
\begin{equation}
 D_N^{(k)}:=\KL(\mathbb P_N\Vert\widehat{\mathbb P}_N^{(k)})
 =\sum_{t=0}^{N-1}
 I_{\mathbb P_N}\bigl(Y_t;Y_{0:t-k-1}\mid Y_{\max(0,t-k):t-1}\bigr),
\label{eq:rg-memory-information-defect}
\end{equation}
where an empty discarded history contributes zero. These are established
finite-law statements, including uniqueness as a path law rather than of
irrelevant conditional completions. \status{ESTABLISHED}

\paragraph{Proof.}
At \(t=0\) the empty context has mass one. Suppose the retained context
has its \(\mathbb P_N\) marginal. Multiplying by its supported row gives
the correct context-plus-symbol marginal, and marginalizing its oldest
symbol gives the next retained context marginal. This induction proves that
no positive constructed prefix enters a null context, and successive row
sums prove normalization. Every positive source word uses positive factors.
For a competing law with rows \(S_t\), expansion of the finite logarithm
products and regrouping by contexts gives
\begin{equation}
 \KL(\mathbb P_N\Vert S)-D_N^{(k)}
 =\sum_{t,c:\mathbb P_N(c_t=c)>0}\mathbb P_N(c_t=c)
     \KL(R_t^{(k)}(\cdot\mid c)\Vert S_t(\cdot\mid c)).
\label{eq:rg-memory-excess-risk}
\end{equation}
If the competing law excludes a positive source word, its KL is infinite
and the inequality is immediate; otherwise all sums are finite. Equality
forces equality of every occupied row. Recursive reachability then forces
the same path law. Finally subtract the exact source chain-rule
logarithm from the projected product; its expected terms are the difference
between entropy conditioned on the retained context and entropy conditioned
on the full past. This is precisely the conditional mutual information in
\cref{eq:rg-memory-information-defect}. \(\square\)

Increasing \(k\) decreases \(D_N^{(k)}\), and \(k\geq N-1\) gives
exact recovery, by conditioning and the displayed information identity.
The projection is idempotent because it retains its own defining
context-plus-symbol marginals. Restriction to a shorter prefix commutes
with the construction, clipping the order only at its available history.
These are established finite consequences, not determination of a finite
memory order for an infinite process. \status{ESTABLISHED}

Forward error controls bounded path observables: for finite probability laws
\(P,Q\), \(2\operatorname{TV}(P,Q)^2\leq\KL(P\Vert Q)\).
To see this, group the KL on \(\{P\geq Q\}\) and its complement; the
remaining within-group divergences are nonnegative. The resulting binary
KL has second derivative \(1/[p(1-p)]\geq4\), value zero, and first
derivative zero at \(p=q\), so it is at least \(2(p-q)^2\).
Boundary cases follow by continuity or infinite KL. This establishes the
stated control, but does not control reverse KL or every recognition-weighted
VFE difference. \status{ESTABLISHED}

For example, the stationary cycle labeled \((0,0,1)\) has
\(\mathbb P_3(000)=0\) but
\(\widehat{\mathbb P}_3^{(1)}(000)=1/6\), so the reverse KL is infinite.
More generally, for separately varied \(\rho\) supported on a positive
source support, the auxiliary VFE change is
\begin{equation}
 \KL(\rho\Vert\widehat{\mathbb P}_N^{(k)})
   -\KL(\rho\Vert\mathbb P_N)
 =\E_\rho\log\frac{\mathbb P_N}{\widehat{\mathbb P}_N^{(k)}}.
\label{eq:rg-memory-recognition-vfe-difference}
\end{equation}
Subtraction of the two finite log ratios proves this established identity.
Its worst-case bound over such \(\rho\) uses the extrema of the supported
log ratio, rather than its \(\mathbb P_N\)-average.
\status{ESTABLISHED}

\section{Stationary entropy, duration, and effective brackets}
\label{sec:rg-memory-stationary}

Now assume \(\nu=\pi\), with \(\pi K=\pi\), and restrict the hidden
set to its positive stationary support. That support is closed: stationarity
and nonnegativity force \(K_{ij}=0\) if \(\pi_i>0=\pi_j\).
Stationary initialization is an additional hypothesis, excluding arbitrary
initial preparations from the rate identities below. \status{HYPOTHESIS}

Set \(H_N=H(Y_0,\ldots,Y_{N-1})\), \(H_0=0\), and
\begin{equation}
 h_j=H(Y_j\mid Y_0,\ldots,Y_{j-1}),\qquad
 h_\infty=\lim_{j\to\infty}h_j,\qquad r_k=h_k-h_\infty.
\label{eq:rg-memory-entropy-notation}
\end{equation}
For \(j=0\) the conditional entropy means \(H(Y_0)\). All entropies
use natural logarithms. The existence of the limit is proved next rather
than added to these definitions. \status{DEFINITION}

\theoremheading{Stationary projection and the finite hidden information budget}{thm:rg-memory-duration}
The stationary projection is normalized, stationary, and prefix-consistent;
it agrees with the source on every window of length at most \(k+1\).
The sequence \(h_j\) decreases to \(h_\infty\), and
\begin{equation}
\begin{aligned}
 D_N^{(k)}&=\sum_{t=k+1}^{N-1}(h_k-h_t),
 &\lim_{N\to\infty}\frac{D_N^{(k)}}N&=r_k,\\
 \sum_{j=0}^\infty r_j&\leq H(Z_0)=:C_Z,
 &r_k&\leq\frac{C_Z}{k+1},\\
 0\leq r_k-\frac{D_N^{(k)}}N&\leq\frac{C_Z}{N},
 &D_N^{(k)}&\leq(N-k-1)_+r_k.
\end{aligned}
\label{eq:rg-memory-duration-identities}
\end{equation}
An empty sum is zero. These are established bounds for any finite stationary
hidden model, without irreducibility or a forgetting assumption.
\status{ESTABLISHED}

\paragraph{Proof.}
For \(k\geq1\), the positive \(k\)-context chain shifts
\(c\) to its last \(k\) symbols after appending \(a\), with row
\(\mathbb P_{k+1}(ca)/\mathbb P_k(c)\). Its incoming mass at a context
\(d\) is the corresponding marginal of \(\mathbb P_{k+1}\), equal
to \(\mathbb P_k(d)\) by stationarity. This proves stationarity and
exact \((k+1)\)-windows; shorter prefixes follow by marginalization.
For \(k=0\), the law is the stationary independent product of the
one-symbol marginal.

Conditioning and stationarity give \(0\leq h_{j+1}\leq h_j\), so the
limit exists. The entropy chain rule gives \(H_N=\sum_{j<N}h_j\),
whose average tends to \(h_\infty\). For \(N>k\), the projected
cross-entropy is \(H_k+(N-k)h_k\), giving the exact KL sum and its rate;
shorter horizons are exact. To bound total excess, take two adjacent finite
observed blocks of lengths \(b,n\). Their past and future are conditionally
independent given the intervening hidden state, so their mutual information
is at most \(C_Z\). Expanding it by stationary entropies gives
\(H_n-\sum_{j=0}^{n-1}h_{b+j}\leq C_Z\). With \(n\) fixed let
\(b\) increase, obtaining \(\sum_{j<n}r_j\leq C_Z\). Increasing
\(n\) proves the nonnegative series bound, and monotonicity gives
\((k+1)r_k\leq C_Z\). Finally, for \(N>k\),
\(Nr_k-D_N^{(k)}=kr_k+\sum_{t=k}^{N-1}r_t\leq C_Z\), while for
\(N\leq k\), \(Nr_k\leq C_Z\). Bounding each nonzero KL summand
by \(r_k\) gives the last inequality. \(\square\)

\theoremheading{Finite hidden-boundary brackets converge to the entropy rate}{thm:rg-memory-effective-brackets}
For \(m\geq1\), define
\begin{equation}
 g_m=H(Y_m\mid Z_0,Y_1,\ldots,Y_{m-1}),\qquad
 d_m=h_{m-1}-g_m.
\label{eq:rg-memory-hidden-bracket}
\end{equation}
The future symbol \(Y_1\) follows one transition from \(Z_0\).
Then
\begin{equation}
 g_m\uparrow h_\infty,\qquad h_{m-1}\downarrow h_\infty,
 \qquad 0\leq d_m\leq C_Z/m,
 \qquad
 \max(0,h_k-h_{m-1})\leq r_k\leq h_k-g_m.
\label{eq:rg-memory-effective-bounds}
\end{equation}
These are established finite-sum bounds and convergent brackets, including
for reducible stationary hidden models. \status{ESTABLISHED}

\paragraph{Proof.}
Adding \(Z_1\) to the conditioning and using the Markov factorization
and \(Y_1=f(Z_1)\) gives
\(g_{m+1}\geq H(Y_{m+1}\mid Z_1,Y_2,\ldots,Y_m)=g_m\).
Stationarity yields
\(d_m=I(Z_0;Y_m\mid Y_1,\ldots,Y_{m-1})\geq0\).
The already proved monotonicity of \(h\), together with that of \(g\),
makes \(d_m\) nonincreasing. The mutual-information chain rule gives
\(\sum_{m=1}^M d_m=I(Z_0;Y_1,\ldots,Y_M)\leq C_Z\), whence
\(d_m\leq C_Z/m\to0\). Therefore both bracket endpoints have the
same limit. Subtracting their inequalities from \(h_k\) proves the rate
bracket. Every entropy here is a finite sum of normalized path probabilities
computed from \(K\); no infinite-past conditional is needed.
\(\square\)

For rational finite models, the brackets give a finite procedure for any
prescribed positive entropy-rate accuracy. Choose
\(m\geq2C_+/\varepsilon\), where \(C_+\geq C_Z\), and enclose each
endpoint to width at most \(\varepsilon/4\). The enclosing bracket then
has width at most \(C_Z/m+\varepsilon/2\leq\varepsilon\).
Such logarithm enclosures exist from the convergent series
\(\log u=2\sum_{j\geq0}s^{2j+1}/(2j+1)\),
\(s=(u-1)/(u+1)\), after reducing a positive rational to
\(2^a u\), \(1\leq u<2\); the positive tail is bounded by a geometric
series. Only finitely many positive ratios occur at a fixed depth. This
establishes finite stopping for accuracy, but not efficient enumeration or
finite resolution of an exact equality to a requested threshold.
\status{ESTABLISHED}

\section{A constructive logarithmic memory schedule under flag hypotheses}
\label{sec:rg-memory-flags}

The following stronger rate uses additional support information. Put
\(T^a_{ij}=K_{ij}\mathbf1_{f(j)=a}\). Assume that for every hidden
target \(j\), a word \(w_j\) of one common positive length \(b\) is
supplied such that, with \(B^w=T^{w_1}\cdots T^{w_b}\),
\begin{equation}
 S_j=\{i:(B^{w_j}\mathbf1)_i>0\}\ne\varnothing,
 \qquad B^{w_j}_{ij}>0\quad(i\in S_j).
\label{eq:rg-memory-flag-hypothesis}
\end{equation}
Thus every origin capable of a selected word can follow that word to its
selected endpoint. One merging pair is insufficient; the condition covers
all generating origins. Positive \(\pi\), fixed \(K\), and these
common flag words are the adopted hypotheses. \status{HYPOTHESIS}

Let \(A=|f(\mathsf S)|\), \(\lambda=\min\{K_{ij}:K_{ij}>0\}\),
and define positive constants
\begin{equation}
\begin{aligned}
 q&=\lambda^b,\qquad \theta=q\min_j\pi(S_j),\\
 a&=\theta^2/2,\qquad c=\theta q/2,\\
 \gamma&=\min(a,c/2),\qquad C_R=2/b+2(A-1).
\end{aligned}
\label{eq:rg-memory-flag-constants}
\end{equation}
These are declared conservative constants; they are not fitted decay rates.
\status{DEFINITION}

\theoremheading{Explicit exponential excess and logarithmic path memory}{thm:rg-memory-logarithmic-schedule}
Under the flag hypotheses, for \(k\geq b-1\),
\begin{equation}
 r_k\leq C_Re^{-\gamma t(k)},\qquad
 t(k)=\left\lfloor\frac{k+1}{b}\right\rfloor-1.
\label{eq:rg-memory-exponential-rate}
\end{equation}
Let \(s_N\) be the least nonnegative integer with
\(C_R2^{-s_N}\leq2/N^3\), and put
\begin{equation}
 k_N=\min\left(N-1,
 b\left(\left\lceil s_N/\gamma\right\rceil+1\right)-1\right).
\label{eq:rg-memory-logarithmic-order}
\end{equation}
For this fixed flagged model, \(k_N=O(\log N)\) and
\begin{equation}
 D_N^{(k_N)}\leq2/N^2,\qquad
 \operatorname{TV}(\mathbb P_N,\widehat{\mathbb P}_N^{(k_N)})\leq1/N.
\label{eq:rg-memory-logarithmic-control}
\end{equation}
This is an established sufficient schedule; its constants may be arbitrarily
poor across different finite models. \status{ESTABLISHED}

\paragraph{Proof.}
Group \(b\) transitions into a block, whose edge-emitted alphabet has size
\(L\leq A^b\). For any supported suffix \(v\), choose a hidden state
\(j(v)\) maximizing its suffix likelihood \(P_j(v)\). Stationarity
and the flag condition give the conditional probability of its preceding
flag as
\begin{equation}
 \frac{\sum_{i,j}\pi_iB^{w_{j(v)}}_{ij}P_j(v)}
      {\sum_j\pi_jP_j(v)}
 \geq\sum_i\pi_iB^{w_{j(v)}}_{i,j(v)}
 \geq q\pi(S_{j(v)})\geq\theta.
\label{eq:rg-memory-flag-occurrence}
\end{equation}
The positive maximum bounds the denominator and cancels from the restricted
numerator. Reveal a stationary word of \(t\) blocks from right to left,
and count these selected flags. Their conditional success probabilities
are at least \(\theta\). For indicators \(I_s\), repeated conditional
expectation bounds
\(\E\exp(-u\sum_sI_s)\) by
\((1-\theta+\theta e^{-u})^t\). The centered Bernoulli log transform
has value and derivative zero at zero and second derivative at most \(1/4\),
so it is at most \(u^2/8\). Exponential Markov inequality with
\(u=2\theta\) proves that fewer than \(\theta t/2\) flags has
probability at most \(e^{-at}\).

Condition a hidden path on one complete supported word. At a selected flag,
with subsequent suffix \(v\), its bridge transition is
\(B^{w}_{ij}P_j(v)/\sum_\ell B^w_{i\ell}P_\ell(v)\).
Every reachable row assigns at least \(q\) to \(j(v)\): bound the
suffix likelihood denominator by its maximum and use a positive length-\(b\)
path, whose probability is at least \(\lambda^b\). Couple two such
bridges with one shared Bernoulli \(q\) coin at each counted flag and
identical transitions after meeting. Their terminal disagreement probability
on a good word is at most
\((1-q)^{\theta t/2}\leq e^{-ct}\); if \(q=1\), meeting is certain
at a counted flag, without a residual row. The next-block channel contracts
this total-variation bound, as do mixtures over initial hidden states.

For laws at total variation \(d\) on an \(L\)-symbol alphabet, a
maximal coupling and its disagreement indicator give
\(|H(u)-H(v)|\leq d\log L+2\sqrt d\). Indeed conditional entropy is
at most binary disagreement entropy plus \(d\log L\), and
\(-d\log d-(1-d)\log(1-d)\leq2\sqrt d\) on \([0,1]\).
Let \(\beta_t\) be conditional block entropy after \(t\) blocks minus
the block entropy rate. Conditioning additionally on the initial hidden
boundary gives a lower bound on that rate: condition a longer stationary
past on that hidden state, use the Markov factorization to remove the remote
past, and then let its length increase. Thus \(\beta_t\) is bounded
by the conditional information between that boundary and the next block.
On good words the entropy-continuity bound applies to each origin-conditioned
prediction and its mixture; on bad words this information is at most
\(\log L\). Consequently, also at \(t=0\) by the trivial entropy bound,
\begin{equation}
 \beta_t\leq2e^{-ct/2}+\log L(e^{-ct}+e^{-at}).
\label{eq:rg-memory-block-excess}
\end{equation}
The entropy chain rule identifies the block rate as \(bh_\infty\), and
\(\beta_t=\sum_{j=bt}^{b(t+1)-1}r_j\geq br_{b(t+1)-1}\).
For the displayed \(t(k)\), the last index is at most \(k\), so
monotonicity, \(\log L\leq b\log A\leq b(A-1)\), and the definition
of \(\gamma\) prove \cref{eq:rg-memory-exponential-rate}.

If the cap \(N-1\) is active, the path law is exact. Otherwise
\(\gamma t(k_N)\geq s_N\) and \(e^{-s_N}\leq2^{-s_N}\), so
\(r_{k_N}\leq2/N^3\). Apply \(D_N^{(k)}\leq Nr_k\) and the proved
finite Pinsker inequality. Finally successive halving finds
\(s_N=O(\log N)\); multiplication by the fixed constants \(b,1/\gamma\)
and the cap preserve this order. This proves the entire finite schedule,
not merely a limiting exponent. \(\square\)

The theorem does not imply uniform contraction on every conditioned word.
For a positive word, any preparation-specific posterior is a mixture of the
normalized nonzero rows of its masked hidden matrix; the supremum distance
between preparations is exactly those rows' diameter, attained at point-mass
starts. Two disjoint killed supports with positive holding entries retain
disjoint posteriors along their common constant-symbol words. This is an
established obstruction to arbitrary-preparation forgetting when that support
configuration exists, compatible with the stationary good-word theorem.
\status{ESTABLISHED}

A direct representation of the projected law needs only its requested
context row,
\begin{equation}
 R_t^{(k)}(a\mid c)=\frac{\mathbb P_{\ell+1}(ca)}{\mathbb P_\ell(c)},
 \qquad \ell=\min(k,t),\quad\mathbb P_\ell(c)>0.
\label{eq:rg-memory-requested-context}
\end{equation}
The exact hidden filter computes these finite word masses, starting from
\(\pi\) with an initial observation mask and transitions between later
symbols. The empty-context row is the stationary observation marginal.
The reachability proof of \cref{thm:rg-memory-projection} proves that
multiplying only requested rows defines the same normalized family; no
complete alphabet-power table is mathematically required. Null direct
queries remain undefined. This is an established representation identity,
with no stronger rate or minimum-order assertion. \status{ESTABLISHED}

Coherent hidden and observed bijections carry all matrices, contexts, flag
words, and probability laws together, preserving these results by reindexing
finite sums. Interpreting such a bijection as a gauge change additionally
requires the separately supplied full observation and parent squares.
Neither this invariance nor memory approximation inserts gauge curvature,
Fisher curvature, or base curvature into the original fixed generative VFE.
These conclusions and exclusions are established consequences of the typed
probability construction. \status{ESTABLISHED}

A useful minimum-memory criterion, uniform constants near a critical limit,
geometry-changing parent transitions, unrestricted nested lineage, and a
physical interpretation of the inference clock remain open. The logarithmic
schedule resolves one fixed finite flagged class; it does not settle those
additional quantifiers. \status{OPEN}

\Cref{ch:finite-depth-formation} gives a complementary finite-depth
construction with actual successive parent activation. Every activation
pushes the complete current law, including all source assignments and loop
marks; subsequent parent learning uses that parent's own evolving law.
Conditional entropy, continuous-trajectory displacement, and complete
assignment-history disagreement receive separate estimates. This
construction neither replaces the retained hidden-history law above nor
turns the finite-memory projection into an exact Markov closure for
arbitrary preparations. \status{ESTABLISHED}
